\section{Threats to Validity}
\label{sec:ttv}
\begin{description}[leftmargin=0.3cm]
\item[Construct Validity.] The psychological signals are proxies inferred from language and reactions, not measurements of psychological states; we therefore read them as \textit{signals} and never as diagnoses. The emotion classifier reaches a micro-F1 of 0.62 on GitHub comments, which improves over existing tools but leaves room for misclassification, especially for fear and anger; moreover, anger includes mild frustration, as in the taxonomy of Shaver \etal~\cite{shaver1987emotion}, so \tbd{[we repeated the analyses considering only explicit anger, \ie with a probability above 0.5]}. We mitigated this threat by validating the classifier on two independent GitHub gold standards and by analyzing relative differences between developers and over time rather than absolute rates. Belonging and gratitude rely on dictionaries, which ignore context, \eg ironic thanks, and apply to English messages only. Leaving GitHub is defined through public contributions, so developers who moved to private repositories may be misclassified; we mitigated this threat by also counting the private contributions that users choose to disclose. Finally, the 12-month threshold for abandonment is a convention~\cite{calefato2022will}, and \tbd{[we assessed the sensitivity of our results to thresholds of 6 and 18 months]}.

\item[Internal Validity.] Our design is observational, so the associations we report are not causal. Reverse causality is plausible: a developer who is already planning to leave may write less with \textit{we} because they disengage, rather than disengaging because their belonging decreases. We mitigated this threat by lagging all covariates, by controlling for activity, role, and project decline, and by entering developers in the risk set without using their future activity. Confounders such as personal events or job changes~\cite{miller2019people} remain unobserved.

\item[External Validity.] \tbd{[Placeholder: generalizability of the selected projects, \eg size, domains, and languages.]} The study is limited to GitHub and to English communication, and the emotion classifier was trained on comments from four large projects and on Stack Overflow posts, which may not represent all communities.

\item[Conclusion Validity.] Developers contribute to several months and, sometimes, to several projects, so observations are not independent. We addressed this threat through cluster-robust standard errors by person, random intercepts for developers, project fixed effects, and stratification by project. We tested the proportional hazards assumption and checked the number of events per covariate to avoid overfitting, and \tbd{[we corrected for multiple comparisons through the Benjamini-Hochberg procedure]}.
\end{description}
